Para celulas geometricamente semelhantes,
\[
\text{area} \sim L^2, \qquad \text{massa} \sim \text{volume} \sim L^3.
\]
Assim,
\[
\frac{\text{area}}{\text{massa}} \sim \frac{1}{L}.
\]
Se a celula se divide quando o volume atinge $1.5$ vezes o valor original, entao o comprimento cresce por
\[
L \mapsto 1.5^{1/3}L.
\]
Logo, a razao area/massa passa a ser multiplicada por
\[
1.5^{-1/3},
\]
isto e, ela diminui. Portanto, esse comportamento nao preserva a proporcao entre superficie e massa, e o pepino nao pode ser descrito como normal segundo o criterio dado.

